% Zdroj: PDF priklady-podzim-2022.pdf, fyzická str. 1-2. Značka: společné okruhy. Zařazeno: I3.
\section{Programování}
\szzotazka{podzim 2022}{1}{Programování (objektové rozhraní pro grafy)}

\noindent\textbf{Zadání.}\par
Pro tuto otázku si vyberte libovolný mainstreamový objektově orientovaný staticky
typovaný jazyk (C++, C\# nebo Java) a svůj výběr napište na začátek řešení.

\begin{enumerate}
\item Naimplementujte (nadeklarujte) vhodné objektové rozhraní (pomocí \textit{interfaces}
  nebo abstraktních tříd) pro neorientované grafy $G = (V, E)$ ve své nejběžnější podobě
  (hrany spojují právě dva vrcholy, mezi dvěma vrcholy vede nejvýš jedna hrana). Motivací
  je, aby toto rozhraní mohly implementovat různé reprezentace grafů (seznam sousedů,
  matice sousednosti, \dots), zatímco programátor grafových algoritmů (BFS, hledání
  minimální kostry, \dots) může napsat svůj kód pouze jednou a bude fungovat se všemi
  reprezentacemi.

  Rozhraní pokrývá 3 entity (vrchol, hranu a graf): vrchol nabídne přístup ke všem
  incidentním hranám, hrana ke svým dvěma vrcholům, graf k seznamu všech vrcholů a všech
  hran. Graf je pouze pro čtení (bez modifikací).

  Každý vrchol a každá hrana má příznak (\textit{tag}) -- u vrcholu např. zda byl
  navštíven, u hrany např. váhu. Tagy mohou mít různé datové typy, ale pro danou instanci
  grafu mají všechny vrcholy stejný typ tagu a všechny hrany taktéž. Rozhraní má být
  generické (uživatel zvolí typy tagů vrcholů a hran) a nabídne pro každý vrchol/hranu
  prostředky tag přečíst a nastavit.

\item Uvažme implementaci, kde tagy vrcholů jsou \texttt{int} (index komponenty
  souvislosti, od~0) a tagy hran jsou \texttt{float} (délky hran, na počátku
  inicializované). Napište tělo funkce využívající vaše rozhraní, která dostane graf a
  reálné číslo $r$ a provede rozklad grafu na komponenty souvislosti tak, že každý vrchol
  bude mít index své komponenty uložený jako tag (konkrétní hodnoty indexů nejsou
  podstatné, jen že vrcholy téže komponenty mají stejný tag). Hrany s délkou (tagem) větší
  než $r$ se neberou v potaz (jako by neexistovaly).
\end{enumerate}

\medskip
\noindent\textbf{Náčrt řešení.}\par
\textit{(V~PDF bez náčrtu řešení; náčrt doplněn k~výkladu okruhu.)}

\medskip
Zvolme \textbf{Javu}. Generika nesou typy tagů: \texttt{VTag} pro vrcholy, \texttt{ETag}
pro hrany.

\begin{enumerate}
\item Tři spolupracující generická rozhraní (pouze deklarace, bez těl):

\begin{lstlisting}[language=Java]
interface IVertex<VTag, ETag> {
    Iterable<IEdge<VTag, ETag>> incidentEdges(); // hrany incidentni s vrcholem
    VTag getTag();
    void setTag(VTag tag);
}

interface IEdge<VTag, ETag> {
    IVertex<VTag, ETag> endpoint1();             // jeden z dvou koncovych vrcholu
    IVertex<VTag, ETag> endpoint2();             // druhy koncovy vrchol
    ETag getTag();
    void setTag(ETag tag);
}

interface IGraph<VTag, ETag> {                   // graf je pouze pro cteni
    Iterable<IVertex<VTag, ETag>> vertices();
    Iterable<IEdge<VTag, ETag>> edges();
}
\end{lstlisting}

  Díky generikám je rozhraní nezávislé na konkrétním typu tagu i na vnitřní reprezentaci
  grafu (seznam sousedů, matice sousednosti, \dots). Algoritmus se proti rozhraní napíše
  jednou a běží nad libovolnou implementací.

\item BFS z každého dosud neoznačeného vrcholu; šíříme se jen po hranách s tagem
  $\le r$, ostatní (delší) hrany ignorujeme, jako by neexistovaly. Všem vrcholům dosaženým
  v jednom průchodu nastavíme stejný index komponenty, čítač pak zvýšíme.

\begin{lstlisting}[language=Java]
static void connectedComponents(IGraph<Integer, Float> g, float r) {
    final int UNVISITED = -1;
    for (IVertex<Integer, Float> v : g.vertices()) {
        v.setTag(UNVISITED);                     // inicializace: nikdo neoznacen
    }

    int component = 0;
    for (IVertex<Integer, Float> start : g.vertices()) {
        if (start.getTag() != UNVISITED) {
            continue;                            // jiz prirazen do komponenty
        }
        // BFS z vrcholu start, vse dosazene dostane index "component"
        Queue<IVertex<Integer, Float>> queue = new ArrayDeque<>();
        start.setTag(component);
        queue.add(start);
        while (!queue.isEmpty()) {
            IVertex<Integer, Float> u = queue.remove();
            for (IEdge<Integer, Float> e : u.incidentEdges()) {
                if (e.getTag() > r) {
                    continue;                    // prilis dlouha hrana, ignoruj
                }
                // urci druhy konec hrany (graf je neorientovany)
                IVertex<Integer, Float> w =
                    (e.endpoint1() == u) ? e.endpoint2() : e.endpoint1();
                if (w.getTag() == UNVISITED) {
                    w.setTag(component);
                    queue.add(w);
                }
            }
        }
        component++;                             // dalsi komponenta
    }
}
\end{lstlisting}

  Hrany s tagem $> r$ se v BFS přeskočí, takže rozklad odpovídá podgrafu jen s hranami
  délky $\le r$. Protože pracujeme výhradně přes rozhraní (\texttt{vertices()},
  \texttt{incidentEdges()}, \texttt{getTag/setTag}), je tělo funkce stejné pro každou
  reprezentaci grafu.
\end{enumerate}

\medskip
\begin{csalt}[generická rozhraní grafu]
Stejný návrh je v~C\# přímočařejší: rozhraní mají vlastní klíčové slovo \texttt{interface},
tag je \emph{vlastnost} s~\texttt{get}/\texttt{set} místo dvojice metod a~kolekce se vracejí jako
\texttt{IEnumerable} (procházené přes \texttt{foreach}). Tělo BFS zůstává nezávislé na reprezentaci grafu.
\begin{lstlisting}[style=cs]
interface IVertex<VTag, ETag> { VTag Tag { get; set; }
    IEnumerable<IEdge<VTag, ETag>> IncidentEdges { get; } }
interface IEdge<VTag, ETag>   { ETag Tag { get; set; }
    IVertex<VTag, ETag> Endpoint1 { get; }  IVertex<VTag, ETag> Endpoint2 { get; } }
interface IGraph<VTag, ETag>  { IEnumerable<IVertex<VTag, ETag>> Vertices { get; }
    IEnumerable<IEdge<VTag, ETag>> Edges { get; } }        // jen pro cteni
\end{lstlisting}
Proti C++ odpadá virtuální destruktor i~ruční správa životnosti uzlů a~hran (uklízí GC).
\end{csalt}
